\section{Research Procedure and Cross-Round Evidence}
\label{app:research}
The main text follows one repair into the policy. The full procedure in
Figure~\ref{fig:framework} chooses those repairs, allocates tests, and
carries completed evidence into later rounds.

\subsection{Hierarchical Failure Analysis}
\label{sec:analysis-memory}
The evidence package $E_k$ contains policy actions, observations, feedback,
mechanically computed progress events, budgets, termination, and checkpoint
provenance. Local analysis $L$ identifies error lifecycles and recovery
opportunities. Deterministic grouping supplies related trajectories to $G$
for recurring-mechanism analysis. Stage $C$ aggregates component and
capability hypotheses. Stage $P$ assembles the global policy-behavior profile
deterministically, and $X$ checks its support, contradictions, and unresolved
alternatives. Here $P$ denotes the profile stage; the task policy and the
planner's bottleneck choice have separate roles.

A0 uses simple local analysis, A1 structured local hypotheses, A2 the
cross-trajectory hierarchy, and A3 the hierarchy plus historical experience.
A2--A1 and A3--A2 define hierarchy and history comparisons, respectively.
No matched effect estimate for these contrasts is reported in this study.
The reused R5 analysis includes 60 local slots, 44 completed G pairs,
88 C terminals, and 87 X terminals. These are analysis records, not
independent successful repair experiments.

\subsection{Procedural Experience across Rounds}
The round-start snapshot $S_k$ stores starting conditions, actions and
feedback, state changes, a wrong branch and its consequence, repair
applicability, counterexamples, and effect provenance. Policy, Analyzer, and
Planner views expose the parts appropriate to each role. A past repair
can enter a new proposal as an analogy or counterexample; current paired
execution determines the new effect.

Writes produced within the round enter a shadow store. Historical reads
throughout round $k$ use the fixed snapshot $S_k$. Only after round closure
do admitted records become part of $S_{k+1}$, which is first readable in
round $k+1$. R3 PRE receives closed R1/R2 lessons, and R4 PRE receives R1--R3 lessons;
repeated local proposals therefore cannot be explained solely by absent
history access. Matched history and ranking comparisons remain unexecuted.
The read/write arrows in Figure~\ref{fig:framework} depict these
mediated interfaces. Current verdicts reach POST directly as experiment
outcomes. This timing keeps current results distinguishable from earlier
experience.

\subsection{Research Planning before and after a Test}
\label{sec:planning}
PRE receives the failure profile, candidate proposals, history, and available
resources. It selects a principal bottleneck, ranks a bounded portfolio,
and fixes the intervention, endpoints, stopping conditions, and budget.
A principal change can be evaluated at several registered states. The
main training portfolio and the discovery-comparison portfolios retain
separate experiment identities.

POST receives the fixed PRE decision and independent verification outcomes
and costs. It considers alternative explanations, recommends TRAIN or
NO-TRAIN, selects eligible supervision, and records the next research
question. Effect labels remain the verifier's output; candidate acceptance
uses the predeclared rule below.

\subsection{Policy Selection and the Next Round}
\label{app:round-procedure}
Let $A_k\in\{0,1\}$ be the predeclared scaffold-free acceptance result on
$\mathcal D_{\mathrm{select}}$. The retained policy for the next round is
\begin{equation}
\pi_{k+1}=\begin{cases}
\widetilde\pi_{k+1},&A_k=1,\\
\pi_k,&A_k=0.
\end{cases}
\label{eq:promotion}
\end{equation}
For the reported final validation, integrity must pass, and the candidate
is accepted if its selection success count increases, or if success ties
and mean terminal goal-condition completion increases. Otherwise the
parent is retained. Goal-condition evaluation is separate from policy
input and does not alter causal Benefit labels. A NO-TRAIN decision also retains $\pi_k$.
Algorithm~\ref{alg:repair} connects this selection to collection and research.
\begin{algorithm}[H]
\caption{RePair research and policy-update procedure}
\label{alg:repair}
\begin{algorithmic}[1]
\Require Initial policy $\pi_0$, history $S_0$, train-side pools, budgets and acceptance rule
\State $k\gets0$; initialize attempt and cost records
\While{the stopping rule permits another attempt}
  \State Fix $\pi_k$, $S_k$, the interface, and the attempt budget
  \State Collect fresh rollouts and construct evidence $E_k$
  \State $\mathcal C_k\gets$ failure analysis using the permitted history views
  \State $\mathcal P_k\gets$ PRE: rank and register a repair portfolio
  \State Execute paired tests; obtain independent effect labels
  \If{the attempt is protocol-invalid}
    \State Record the incident and cost; apply the registered recovery or stop
    \State \textbf{continue}
  \EndIf
  \State $d_k\gets$ POST: interpret outcomes and recommend TRAIN or NO-TRAIN
  \State Construct $\mathcal D_k^+$ under Eq.~\eqref{eq:eligibility}
  \If{$d_k=\mathrm{TRAIN}$ and $\mathcal D_k^+\ne\varnothing$}
    \State Check native labels; stop/recover on a failed label check
    \State Train $\widetilde\pi_{k+1}$ on the registered dual views
    \State Retain $\pi_{k+1}$ by Eq.~\eqref{eq:promotion}
  \Else
    \State Record NO-TRAIN; $\pi_{k+1}\gets\pi_k$
  \EndIf
  \State Close the round and release eligible history into $S_{k+1}$
  \State Update attempt, patience, and cost records; $k\gets k+1$
\EndWhile
\State \Return retained policy and experiment records
\end{algorithmic}
\end{algorithm}


The campaign permits at most ten scientifically valid attempts and stops
after three consecutive valid attempts without promotion, subject to its
additional integrity and resource stops. Invalid attempts follow the
registered recovery procedure; their costs and human recovery actions
remain in the campaign record. The observed stop is the three-valid-round no-promotion patience limit.
The five-round record includes one earlier user-authorized promotion
exception; it is not represented as an automatic acceptance decision.
Strong research models can run this sequence after initial configuration;
its unattended execution is measured separately from policy improvement.

\subsection{Cross-Round Outcomes}
\label{sec:round-results}

The development campaign records five valid rounds and one invalid
attempt. R1, R3, R4, and R5 yield no training update; R2 contains the
beneficial intervention and the user-authorized promotion exception.
Consequently, the sequence is not evidence of five successful updates
or five rounds of uninterrupted automatic improvement. Engineering
recoveries are documented separately and are not erased by resuming.

The later validations reuse R5 trajectories and Analyzer outputs because
the retained parent identity is unchanged. The guarded validation ends
without a training update. The executable-strategy validation trains one
candidate and completes selection with a rollback decision. A subsequent
closed-history ownership check stops its bookkeeping closeout; this does
not change the already completed episode outcomes or turn the run into
a closed new formal round.

The selection entry initially scheduled five seed slots through a legacy
path. An explicit correction retains the prescribed seed 17, after
episode publication, and stops the additional jobs. Both arms have all
355 seed-17 tasks. The supplementary record discloses this amendment;
extra-seed work is not counted as additional evidence or represented as
having never occurred. Goal fractions are recomputed by replaying the
recorded actions with no new policy calls.

